\section{Methods}
\subsection{Models}
Five components were used: logistic regression (LR), RBF support vector machine (SVM), $k$-nearest neighbors (KNN, $k{=}15$), random forest (RF, 500 trees) and XGBoost \cite{xgboost}, each a pipeline whose imputation and scaling are fitted only on its training data. The \emph{hybrid ensemble} (HE) soft-votes within a linear/kernel/instance group \{LR, SVM, KNN\} and a tree group \{RF, XGB\}, then stacks the two group probabilities with an L2 logistic meta-learner trained on inner 5-fold out-of-fold predictions \cite{wolpert}. The comparator S* is the component with the highest inner out-of-fold AUC on the same training fold, i.e., the single model a practitioner would choose without an ensemble. For robustness we also used soft voting (SV) and logistic stacking (ST) of all five components.

\subsection{Explanations and faithfulness}
Both models are explained on the same patients with permutation SHAP (independent masker, 50 complete-case training rows as background, 500 evaluations, common random numbers per patient). Let $s=\pm1$ be the sign of the predicted class, $F(S)$ the mean prediction after the features in $S$ are replaced by their values in 50 further training rows (\emph{marginal} operator), $[d]$ the set of all $d$ features, $G=s\,(f(x)-F([d]))$ the prediction gap, and $T_k$ the $k$ features with the largest positive signed attributions $s\,\phi_j$, with $K=\min(5, \#\{j: s\phi_j>0\})$. With $D(S)=s\,(f(x)-F(S))/G$, the share of the gap removed by deleting $S$,
\begin{equation}
\mathrm{NDF}(x)=\frac{1}{K}\sum_{k=1}^{K}\Big(D(T_k)-\overline{D}^{\,\mathrm{rand}}_k\Big),
\end{equation}
where $\overline{D}^{\,\mathrm{rand}}_k$ is the mean of $D$ over the first $k$ features of 20 random orderings. NDF is 0 when the top-ranked features remove the prediction no faster than random features and positive when they remove it faster. A \emph{conditional} operator instead takes replacement values from the 20 training rows nearest to $x$ on the kept features.

\emph{Assumptions.} (i) Marginal replacement assumes that features can be changed one at a time; with correlated features it evaluates the model on combinations that do not occur in patients, so NDF then measures faithfulness to the model's behavior under that intervention, not to the data. The conditional operator assumes that 20 nearest neighbors approximate the conditional distribution. Because the two answer different questions \cite{chen}, every faithfulness conclusion below names its operator. (ii) $D$ is a share of the gap and is defined only for $G>0$; it is unstable when $G$ is near zero, so the primary analysis excludes patients with $G<0.05$ for either model. We vary this threshold and also report an exclusion-free version, the absolute deletion faithfulness $\mathrm{ADF}=\mathrm{NDF}\times G$ (prediction drop beyond random deletion, in probability units). (iii) NDF is unbounded: $D$ exceeds 1 when deleting a few features overshoots the all-features-deleted prediction. When $d\le K$ (Mammographic, $d=4$) the $k=d$ term is zero by construction. NDF passed known-answer tests on synthetic models before use (Supplement~S1).

\subsection{Datasets, audit and validation}
Eight binary tasks were used (Table~\ref{tab:data}) \cite{street,antal,detrano,little,uci_repo,uci_cdc,brfss}: three pre-specified datasets (\emph{primary}), four analyzed under a second protocol written before their analysis (\emph{replication}), and a stratified sample of 5{,}000 respondents of the CDC Diabetes Health Indicators survey (\emph{large-sample}). The datasets differ in size, feature types and correlation and come from independent sources; we treat them as separate case studies. Table~\ref{tab:data} gives the audit: label definition, duplicates (removed when they indicate repeated entries, kept when coarse features make identical patients plausible), missing values (imputed inside each training pipeline; explanations used complete cases) and repeated subjects (Parkinsons, folds grouped by person). Each dataset was evaluated with $5\times5$ repeated stratified cross-validation; in every outer fold, 20 complete-case test patients formed the explanation sample.

\subsection{Hypotheses and statistics}
\emph{Confirmatory tests.} H1 (equivalence): $|\Delta\mathrm{AUC}|$ of HE versus S* is below the pre-specified margin of 0.02; tested with two one-sided tests (TOST, 90\% CI within $\pm0.02$) on the 25 fold-level differences, with the corrected repeated cross-validation variance \cite{nadeau}. H2: NDF differs between HE and S* under marginal removal, tested with the same corrected $t$-test and Holm adjustment within each protocol \cite{holm}. The margin was fixed in the protocol; its formal TOST test was added afterwards. Dataset estimates were pooled by DerSimonian--Laird (DL) random effects \cite{dersimonian} with $I^2$ \cite{higgins} and, because eight studies are few, also by the more conservative Hartung--Knapp--Sidik--Jonkman (HKSJ) method \cite{inthout}, with 95\% prediction intervals \cite{riley}.

\subsection{Further analyses}
\emph{Simulation.} 600 patients with ten Gaussian features, an equicorrelated block of four ($\rho\in\{0,0.5,0.9\}$), a linear or non-linear logistic outcome, 3 replicates $\times$ 5 folds; HE versus LR.
\emph{Robustness.} (i) Held-out repetitions 2--5; (ii) KernelSHAP and faithfulness correlation \cite{bhatt} (repetition 1, exploratory); (iii) components tuned by nested grid search (repetition 1); (iv) NDF of each component, SV and ST, and HE versus RF and XGB; (v) exact TreeSHAP; (vi) Monte-Carlo error; (vii) gap thresholds 0.02 and 0.10, ADF, and patients with similar gaps. Analyses (iv)--(vii) and the HKSJ and prediction intervals were added after the primary results were known and are \emph{post hoc}.
\emph{Remedy.} Attributions were summed within clusters of correlated features ($|\rho|\ge0.7$) and NDF computed by deleting whole clusters; found on repetition 1 and tested on repetitions 2--5 against a random-cluster placebo (Supplement S11).
